\scanpage{114}

d'un él. de $\mathcal E_{X\deleted{\unsure{/Y}}}^T$, opérant sur le sous-groupe $\mathcal E_{X/Y}^T$ par autom. intérieurs.

Les notations \inserted{$\mathcal E_X$,} $\mathcal E_{X/Y}$ évoquent intuitivement un $p^t$ de vue qui regarde $\mathcal E_{X\deleted{\unsure{/Y}}}^T$ comme une \uline{extension} de $\mathcal E_Y^T$ (ou d'un ss groupe convenable \unsure{seulement}) par $\mathcal E_{X/Y}^T$, qui sera souvent, pour les besoins de la cause, noté $\pi$ --- il joue le rôle d'un groupe fondamental (géométrique), tandis que $\mathcal E_Y^T$ joue le rôle d'un groupe de Galois --- ou d'un groupe profini mixte.

\margin{et j'hésite encore entre les \uline{notations} $\mathcal E_{X/Y}^T$ et $\pi_{X/Y}^T$; que je vais peut-être employer concurremment.}

Le cas le plus important pour nous sera celui où $Y=\operatorname{Spec}K$, $K$ un corps. \deleted{\unsure{Assez} souvent} Dans ce cas, $\mathcal E_Y\simeq\mathcal E_Y^{\widetilde Y}$

\margin{113}
